\section{From phase kinetics to numerical capillary resistance}
\label{sec:phasekinetics}
\subsection{Exact finite-grid spectrum and conditional expansion}
The leading bulk phase equation for the rest-state projection map is the
established conservative Allen--Cahn form \cite{subhedar2020}:
\begin{equation}\label{eq:leadingphase}
 \partial_t\varphi+\nabla\cdot(\varphi\bm u)
 =D\nabla\cdot[\nabla\varphi-2b\varphi(1-\varphi)\bm n].
\end{equation}
Here $D=1/6$ in the common diffusive units, with $b=\beta/h$ in the formal
smooth projection expansion. The independent logistic reference calibrates
its own finite-grid $b$ from the same measured crossing width.
For a smooth signed-distance field $d$,
\begin{equation}
 \varphi=\frac{1+\tanh(bd)}2,\qquad |\nabla d|=1,
\end{equation}
the nonadvective flux vanishes pointwise in the interface neighborhood.
This known profile-preservation property \cite{mirjalili2020,hao2024}
does not imply vanishing shape relaxation for a discretized map.

We differentiate $P$ at its independently equilibrated planar bridge,
including variation of $\bm n_h$ and the production color-gradient cutoff.
Freezing the normal would spuriously retain leading tangential diffusion.
For reflected transverse mode $q=m\pi/H$, let $T(q)$ denote the exact
Fourier matrix and $\mu(q)$ the real eigenvalue with greatest overlap with
planar translation. The predicted physical decay is
\begin{equation}\label{eq:phaserate}
 \Gamma(q)=-\log\mu(q)/h^2,\qquad k=m\pi/H_p.
\end{equation}
No measured decay rate enters the matrix. Its analytic link summation,
active-mask scope and independently differentiated action are supplied
in the supplement.

The full solver is perturbed with paired amplitudes $\pm.01$ and, at H16,
$\pm.02$ in physical length units. Both interfaces have the same displacement;
surface tension and body force are zero. There are \MajorWaveCount\
complete trajectories: modes 1 at H16/H32/H48, modes 2 and 3 at H16/H32,
and mode 1 at separately prescribed $W_{0.9}=6$ on H16/H32.
The maximum relative rate difference from Eq.~\eqref{eq:phaserate} is
\MajorWavePercent\%, without refitting for mode or width.
Moving the fit-window start from 50\% to 40\% or 60\% of available samples
changes rates by at most $\MajorWindowDifference$ relatively.
The reported windows and actual output times remain in the dataset.

\begin{figure}[htbp]\centering
\includegraphics[width=.98\linewidth]{major-spectrum.pdf}
\caption{Finite-grid phase kinetics. Left: exact rates divided by $k^4$
for three modes and two widths; the $h^2/24$ line is a formal guide.
Right: H16 full-solver late amplitudes (symbols) and exact spectral rates
(lines), normalized at each displayed window start. Only amplitude is
anchored; rates are predicted independently. All 26 signed/amplitude
controls are retained in the data.}
\label{fig:majorspectrum}\end{figure}

For an isolated smooth planar profile, projection weights, inactive tail
cutoff and decaying perturbations, a formal expansion gives
\begin{equation}\label{eq:fourthshape}
 \eta_t=-K_{\mathrm{num}}\eta_{yyyy}+O(h^4),\qquad
 K_{\mathrm{num}}=\frac{Dh^2}{4}=\frac{h^2}{24}.
\end{equation}
The unit-normal correction is essential in this coefficient.
The derivation and separate no-cutoff coefficient-refinement check are in
the supplement. This is a conditional long-wave result for the tested
$D=1/6$, not a universal coefficient for arbitrary interfaces or other
mobility models. Exact finite-grid eigenvalues supply the wave predictions,
including the finite-mode departures from the formal guide.

Cutoff changes require a new stationary-profile check.
At fixed original profiles, some $10^{-10}$ or $10^{-8}$ trial thresholds
produce residuals of order $10^{-11}$--$10^{-10}$; their eigenvalues are
not admitted as equilibrium rate predictions. Eligible H32 baseline and
W=6 controls preserve the measured spectrum. This matrix sensitivity
is not a full-solver test with altered production cutoffs or a theorem
of cutoff independence.

\subsection{Selective fourth-order intervention}
The original matched-width $\lambda=1$ dynamic comparison has
$\overline u_{LB}/\overline u_{FV,0}-1$ equal to 17.12\%, 15.18\% and
13.72\% at H16/H32/H48. A prior shared-projection-stencil ablation narrows
the finest discrepancy to 0.290\%. Because it retains the complete
discrete phase operator, that ablation alone neither identifies a
particular term nor validates the leading continuum equation.

We therefore add only
\begin{equation}\label{eq:intervention}
 \mathcal R_a(\varphi)=-aK_{\mathrm{num}}L_y^2\varphi,\qquad
 a=0,\;0.5,\;1,
\end{equation}
to the independent FV right-hand side.
$L_y$ is the cell-centered reflected Neumann Laplacian.
The coefficient is fixed by Eq.~\eqref{eq:fourthshape}, without using a
front velocity. The same planar profile, leading diffusion, measured width,
Stokes solver, surface force and prescribed hydraulic slip are retained.
The added term conserves phase mass and vanishes exactly on a planar bridge.
Its cosine eigenrate is independently checked as
\begin{equation}
 \Gamma_{4,h}=aK_{\mathrm{num}}
       [4h^{-2}\sin^2(m\pi/2H)]^2.
\end{equation}
This diagnostic is transverse and uses reflected fourth-order conditions;
it is not an orientation-independent replacement for $P$ or a proposed
physical phase law.

\begin{table}[htbp]\centering\small
\caption{Prescribed-slip intervention. Mean speeds are multiplied by $10^5$; $\epsilon_1=100(\overline u_{FV,1}/\overline u_{LB}-1)$. The coefficient of the added operator is theoretical and is not fitted.}\label{tab:majorintervention}
\begin{tabular}{rrrrrr}\toprule
$\lambda$ & H & LBM & FV, a=0 & FV, a=1 & $\epsilon_1$ (\%)\\\midrule
1 & 16 & 1.510667 & 1.289872 & 1.526120 & 1.023\\
1 & 32 & 1.394267 & 1.210531 & 1.400568 & 0.452\\
1 & 48 & 1.344230 & 1.182072 & 1.347929 & 0.275\\
4 & 16 & 2.960897 & 2.809180 & 2.983938 & 0.778\\
4 & 32 & 2.873068 & 2.763845 & 2.881916 & 0.308\\
4 & 48 & 2.838763 & 2.749920 & 2.843750 & 0.176\\
\bottomrule\end{tabular}\end{table}


The $a=1$ reference differs from LBM by \MajorCoarsePercent\%,
\MajorMiddlePercent\% and \MajorFinePercent\%, using
$\overline u_{FV,1}/\overline u_{LB}-1$.
The separately prescribed $\lambda=4$ cases give 0.778\%, 0.308\%
and 0.176\%. Half-strength H16/H32 results lie between the untreated
and unit-strength responses (Fig.~\ref{fig:majorintervention}).
The residuals remain numerical discrepancies; no contact-line validation
or continuum extrapolation follows from their small values.

\begin{figure}[htbp]\centering
\includegraphics[width=.98\linewidth]{major-intervention.pdf}
\caption{Full nonlinear reference intervention at two prescribed hydraulic
slips. LBM, FV and the held-normal-profile reduction use independently
specified parameters. Crosses are reduced-model predictions, not speed fits.
Half strength is additionally tested at slip 1 on H16/H32.
The theoretical fourth-order term accounts for most of the dynamic gap;
the agreement is a numerical mechanism test.}
\label{fig:majorintervention}\end{figure}

At H32, reversed force with half time step changes the corrected mean-speed
magnitude by $7.48\times10^{-8}$ relatively; doubling the positive force
changes corrected mean mobility by $2.33\times10^{-7}$.
Relative mass changes are below $2\times10^{-13}$ for these FV controls.
The original initialization cross-control, signed LBM and duration checks
remain separate in the supplement.

\subsection{Predictive low-force resistance reduction}
To predict the intervention rather than only observe it, hold the
independently calibrated normal profile $p(x)$ and write
\begin{equation}
 \delta\varphi(x,y)=-p'(x)\eta(y).
\end{equation}
The derivative has its appropriate sign at each interface.
Project the phase rate onto the ascending-interface integral.
The body-force Stokes response supplies $g_hF$; linearized capillary force
and the same independent Stokes solver supply $\sigma C_h\eta$.
Neither calculation uses an evolved front, an actual-force audit,
a fitted slip nor a measured front velocity.

The reduced displacement equation is
\begin{equation}\label{eq:reducedphase}
 \eta_t=g_hF+A_a\eta,\qquad
 A_a=\sigma C_h-aK_{\mathrm{num}}L_y^2 .
\end{equation}
A uniformly translating shape satisfies
\begin{equation}\label{eq:steadymobility}
 A_a\eta+g_hF=U_\varphi\bm1,\qquad \bm1^T\eta=0.
\end{equation}
The second equation fixes the displacement gauge.
Since rigid translation is a right null mode, a normalized left null vector
$\ell_a$ gives
\begin{equation}\label{eq:leftmobility}
 \frac{U_\varphi}{F}
 =\frac{\ell_a^Tg_h}{\ell_a^T\bm1},\qquad \ell_a^TA_a=0.
\end{equation}
$U_\varphi$ is the projected phase-displacement speed.
The mean fluid speed is calculated separately from the capillary/body
Stokes field induced by $\eta$; they are not assumed identical.

The capillary matrix uses centered profile perturbations of size $10^{-4}$.
Doubling that differentiation amplitude changes the H32 predictions by
less than $9\times10^{-9}$ relatively. Its small rigid-mode differentiation
residue, below $1.4\times10^{-11}$, is recorded before imposing the known
translation identity. Left-null and bordered-system speeds agree within
$3.1\times10^{-14}$ relatively in the retained matrix cases.
The calculated operators have no positive eigenvalue above roundoff;
no nonlinear stability theorem is inferred.

For the 14 baseline nonlinear-reference comparisons at $\sigma=.005$,
the reduced mean-speed error is at most \MajorReducedBaselinePercent\%.
Across all 22 available slip/grid/strength/surface-tension comparisons it
is at most \MajorReducedAllPercent\%.
The held-profile approximation is explicitly checked, not assumed exact:
its phase-volume shapes differ from the full FV reference by up to
\MajorFVShapePercent\%.
For LBM at $\sigma=.005$, the corresponding unit-strength predicted shapes
differ by at most \MajorShapePercent\% over the six slip/grid cases.
This observable is a row-wise phase-volume integral with its mean removed,
not a fitted contour (Fig.~\ref{fig:majorlocalization}).

The resistance mechanism is visible in Eq.~\eqref{eq:leftmobility}.
At $a=0$, scaling $\sigma C_h$ changes the balancing deformation inversely
with $\sigma$, while its left null vector and reduced mobility are unchanged.
At $a>0$, the ratio $K_{\mathrm{num}}/\sigma$ changes $\ell_a$ and therefore
mobility. Thus an interface-shape relaxation can create surface-tension
sensitivity in the linear creeping-flow response without a change in
hydraulic slip. Wall hydrodynamics mix transverse modes, and symmetric
parabolic driving excites even modes. A single first-mode decay rate or
an unbounded-interface estimate is insufficient to predict this resistance.

\subsection{Surface-tension and wall/bulk controls}
We test $\sigma=.0025$ and .01 at H16/H32, retaining the same production
kernel, width, force and slip. Full LBM and FV strengths 0 and 1 are run
independently; low-tension FV0 cases are extended to time 16000 because
their initial duration left a resolvable late drift. The H32 low-tension
LBM is continued from its exact outgoing donor to that same duration.
Both earlier and later histories remain in the data.

The LBM mean speed changes by about 9.11\% at H16 and 7.05\% at H32 over
this fourfold surface-tension range. The untreated FV response changes
only about 0.05\% and 0.094\%, respectively.
The $a=1$ reference reproduces the direction and most of the magnitude,
with LBM discrepancies 0.661--1.574\% at H16 and 0.298--0.673\% at H32.
The reduced model predicts these reference responses without a new
coefficient or mobility fit (Fig.~\ref{fig:majorsigma}).
Its untreated surface-tension independence is a held-profile linear
prediction; small profile/discretization effects in the full FV solution
are retained rather than declared zero.

\begin{figure}[htbp]\centering
\includegraphics[width=.98\linewidth]{major-sigma.pdf}
\caption{Surface-tension controls at fixed physical width, force and slip.
Crosses are independently constructed reduced predictions. The fourth-order
term predicts the numerical surface-tension sensitivity; the untreated
reduced mobility is independent of tension. Long low-tension controls use
their actual extended output times; all earlier observations are retained.}
\label{fig:majorsigma}\end{figure}

To distinguish wall and bulk sensitivity, prescribe
\begin{equation}
 w_w(y)=e^{-(y/2)^2}+e^{-[(H_p-y)/2]^2},\qquad w_b=1-w_w,
\end{equation}
and replace $L_y^2$ by $L_y[wL_y]$.
The two operators sum to the uniform intervention, conserve phase mass,
annihilate the planar profile and have nonpositive quadratic work.
The two-unit wall window is specified before comparing responses;
it is not a measured contact-line length.

At H48, wall-only, bulk-only and uniform mean speeds are
$1.34333$, $1.20847$ and $1.34793$, respectively, multiplied by $10^{-5}$,
compared with LBM $1.34423\times10^{-5}$ and untreated FV
$1.18207\times10^{-5}$.
H32 gives the same ordering (Fig.~\ref{fig:majorlocalization}).
The principal sensitivity to this operator is therefore in the prescribed
wall neighborhood. Responses are not additive even though the operators
are, so a percentage decomposition of physical dissipation is not inferred.

\begin{figure}[htbp]\centering
\includegraphics[width=.98\linewidth]{major-localization.pdf}
\caption{Location and shape controls. Left: conservative wall/bulk/unit
interventions normalized by the full LBM speed; the dotted line is LBM.
Right: LBM row-wise phase-volume displacement and the unit-strength reduced
prediction, with only the displacement gauge removed. The localization
window is prescribed and the profile reduction is approximate. These
comparisons do not measure molecular wall dissipation.}
\label{fig:majorlocalization}\end{figure}
